\section*{致谢}
感谢 Peter Weisz 对本文的批判性阅读并核验了第 2 节中的若干公式。全部数值模拟都是
在 CERN 的一台专用 PC 机群上完成的。感谢 CERN 管理层提供所需经费，并感谢 CERN
IT 部门的技术支持。

\appendix

\section{记号约定}\label{app:A}

$\SUn$ 的李代数 $\sun$ 可以等同于所有反厄米无迹 $N\times N$ 矩阵构成的线性空间。
取这类矩阵的一组基 $T^a$（$a=1,\ldots,N^2-1$），元素 $X\in\sun$
便为 $X=X^aT^a$，其中分量 $X^a$ 为实数（重复的群指标自动求和）。在对易关系
\begin{equation}
  [T^a,T^b]=f^{abc}T^c
  \label{eq:commutator}
\end{equation}
中，若施加归一化条件
\begin{equation}
  \tr\{T^aT^b\}=-\frac{1}{2}\delta^{ab}
  \label{eq:normalization}
\end{equation}
则结构常数 $f^{abc}$ 为实数且对指标完全反对称。此外，
$f^{acd}f^{bcd}=N\delta^{ab}$。

连续理论中的规范场取值于规范群的李代数。其归一化通常使规范变换与协变导数不显含
规范耦合（但参见 2.2 小节）。$D$ 维中的洛伦兹指标 $\mu,\nu,\ldots$
取值 $0$ 到 $D-1$，成对出现时自动求和。时空度规取为欧几里得度规；特别地，对任意
动量有 $p^2=p_{\mu}p_{\mu}$。

本文考虑的格点理论像通常那样建立在间距为 $a$
的四维超立方格点上。特别地，格点规范场 $U(x,\mu)$ 位于格点的链 $(x,\mu)$
上并取值于规范群。作用在规范场函数 $f(U)$ 上的链微分算符为
\begin{equation}
  \partial^a_{x,\mu}f(U)=
  {\rmd\over\rmd s}f(\rme^{sX}U)\!\left.\vphantom{\frac{r}{s}}\right|_{s=0},
  \hskip1.7em
  X(y,\nu)=\begin{cases}T^a & \text{若 $(y,\nu)=(x,\mu)$,}\\[1ex]
                        0  & \text{其他情形,}\end{cases}
  \label{eq:linkderivatives}
\end{equation}
虽然这些算符依赖于生成元 $T^a$ 的选取，但组合
\begin{equation}
  \partial_{x,\mu}f(U)=T^a\partial^a_{x,\mu}f(U)
  \label{eq:covlinkderivative}
\end{equation}
可以证明是与基无关的。

\section{一圈积分}\label{app:B}

除 ${\cal E}_0$ 之外、对 $g_0^4$ 阶的 $\langle E\rangle$
有贡献的费曼积分都涉及对两个动量 $q$ 与 $r$ 以及至多两个流时间参数的积分。在所有
情形，被积函数都具有形式
\begin{equation}
  \rme^{-sq^2-ur^2-v(q+r)^2}{P(q,r)\over q^2r^2(q+r)^2}
\end{equation}
其中 $P(q,r)$ 是洛伦兹不变多项式，而 $s,u,v$
是流时间参数的线性组合。

如前所述，这些积分原则上可以通过对传播子 $1/q^2$、$1/r^2$ 和 $1/(q+r)^2$
代入费曼参数表示来求值。不过，利用下面的基本积分总能进行更经济的计算：
\begin{align}
  \int_q\rme^{-sq^2}(q^2)^{-\alpha}&={s^{\alpha}\over(4\pi s)^{D/2}}
  {\Gamma(D/2-\alpha)\over\Gamma(D/2)},
  \\
  \int_{q,r}{\rme^{-sq^2-ur^2-v(q+r)^2}\over q^2}&=
  {2\over(4\pi)^D(D-2)(u+v)}(su+uv+vs)^{1-D/2}.
  \label{eq:basicintegralB}
\end{align}
例如，一旦把传播子 $1/r^2$ 换成其费曼参数表示，第二个公式便使积分
\begin{equation}
  \int_{q,r}
  {\rme^{-t(q^2+r^2+(q+r)^2)}\over q^2r^2}=
  {1\over(4\pi)^4(2t)^2}\bigl\{8\ln2-4\ln3+\rmO(\eps)\bigr\}
\end{equation}
可以被迅速求出。

一类特殊情形是形如下式的积分
\begin{equation}
\begin{split}
  \int_0^t\rmd s&\int_{q,r}
  {\rme^{-(t+s)(q^2+r^2)-(t-s)(q+r)^2}\over q^2r^2}\,
  qr\\
  &={1\over(4\pi)^4(2t)^2}\bigl\{\frac{1}{2}-4\ln2+2\ln3+\rmO(\eps)\bigr\},
\end{split}
\end{equation}
其被积函数正比于 $qr$。注意到
\begin{equation}
  qr=\frac{1}{2}\left\{(q+r)^2-q^2-r^2\right\},
\end{equation}
这个因子可以换成对流时间参数的求导。多数情况下，这使其中一个参数积分得以立刻完
成，从而给出不含 $qr$ 因子的较简单积分。


\section{威尔逊流的数值积分}\label{app:C}

在有穷格点上，所有规范场的空间 ${\cal G}$ 是规范群的有穷次幂，因而本身就是一个李
群。相应的李代数 $\mathfrak{g}$ 与取值于规范群之李代数的全部链场构成的线性空间重
合。从这个抽象观点看，流方程~\eqref{eq:wflow} 是如下形式的常微分方程
\begin{equation}
   \dot{V_t}=Z(V_t)V_t,
   \label{eq:flowODE}
\end{equation}
其中 $V_t\in{\cal G}$，$Z(V_t)\in\mathfrak{g}$。

本附录描述的 Runge--Kutta 格式从 $t=0$
处的初始组态出发，递推地求得流方程在时刻 $t=n\eps$（$n=1,2,3,\ldots$）
处的解。从时间 $t$ 到 $t+\eps$ 的积分规则是
\begin{align}
  W_0&=V_t,
  \nonumber\\
  W_1&=\exp\bigl\{\frac{1}{4}Z_0\bigr\}W_0,
  \nonumber\\
  W_2&=\exp\bigl\{\frac{8}{9}Z_1-\frac{17}{36}Z_0\bigr\}W_1,
  \nonumber\\
  V_{t+\eps}&=\exp\bigl\{\frac{3}{4}Z_2
                        -\frac{8}{9}Z_1+\frac{17}{36}Z_0\bigr\}W_2,
  \label{eq:RKstep}
\end{align}
其中
\begin{equation}
  Z_i=\eps Z(W_i),\qquad i=0,1,2.
\end{equation}
注意该规则是完全显式的。此外，由于规范场可以从一个方程到下一个方程被覆写，且
$Z_0$ 可以被 $\frac{8}{9}Z_1-\frac{17}{36}Z_0$
覆写，所需的中间存储空间只有后一类场中的一个。

直接计算表明，积分格式~\eqref{eq:RKstep} 每步精确到 $\eps^4$
阶误差；积分到指定流时间的总误差按 $\eps^3$ 标度。经验上发现，该积分在正流
时间方向上是数值稳定的：若 $\eps=0.01$，链变量中的积分误差为 $10^{-6}$ 量级。



\begin{thebibliography}{20}

\bibitem{Wilson}
K. G. Wilson, \textit{Confinement of quarks},
Phys. Rev. D10 (1974) 2445

\bibitem{ArnoldI}
V. I. Arnold,
\textit{Ordinary differential equations}, 3rd ed.\ (Springer-Verlag, Berlin, 2008)

\bibitem{Stout}
C. Morningstar, M. Peardon,
\textit{Analytic smearing of SU(3) link variables in lattice QCD},
Phys. Rev. D69 (2004) 054501

\bibitem{TrivMaps}
M. L\"uscher,
\textit{Trivializing maps, the Wilson flow and the HMC algorithm},
Commun. Math. Phys. 293 (2010) 899

\bibitem{BardeenEtAl}
W. A. Bardeen, A. J. Buras, D. W. Duke, T. Muta,
\textit{Deep inelastic scattering beyond the leading order
in asymptotically free gauge theories},
Phys. Rev. D18 (1978) 3998

\bibitem{SommerScale}
R. Sommer,
\textit{A new way to set the energy scale in lattice gauge theories
and its applications to the static force and $\alpha_s$ in
SU(2) Yang--Mills theory},
Nucl. Phys. B411 (1994) 839

\bibitem{GuagnelliEtAl}
M. Guagnelli, R. Sommer, H. Wittig (ALPHA collab.),
\textit{Precision computation of a low-energy reference scale in
quenched lattice QCD},
Nucl. Phys. B535 (1998) 389

\bibitem{DelDebbioTauQ}
L. Del Debbio, H. Panagopoulos, E. Vicari,
\textit{$\theta$-dependence of SU(N) gauge theories},
JHEP 0208 (2002) 044

\bibitem{StefanConference}
S. Schaefer, R. Sommer, F. Virotta,
\textit{Investigating the critical slowing down of QCD simulations},
PoS (LAT2009) 032

\bibitem{HairerEtAl}
E. Hairer, C. Lubich, G. Wanner,
\textit{Geometric numerical integration:
Structure-preserving algorithms for ordinary differential equations},
2nd ed.\ (Springer, Berlin, 2006)

\bibitem{LambdaParm}
S. Capitani, M. L\"uscher, R. Sommer, H. Wittig (ALPHA collab.),
\textit{Non-perturbative quark mass renormalization in quenched lattice QCD},
Nucl. Phys. B544 (1999) 669

\bibitem{RitbergenEtAl}
T. van Ritbergen, J. A. M. Vermaseren, S. A. Larin,
\textit{The four-loop $\beta$-function in Quantum Chromodynamics},
Phys. Lett. B400 (1997) 379

\bibitem{TopCharge}
M. L\"uscher,
\textit{Topology of lattice gauge fields},
Commun. Math. Phys. 85 (1982) 39

\bibitem{PhillipsStone}
A. Phillips, D. Stone, \textit{Lattice gauge fields,
principal bundles and the calculation of the topological charge},
Commun. Math. Phys. 103 (1986) 599

\bibitem{IndThm}
P. Hasenfratz, V. Laliena, F. Niedermayer,
\textit{The index theorem in QCD with a finite cutoff},
Phys. Lett. B427 (1998) 125

\bibitem{ChiGW}
L. Del Debbio, L. Giusti, C. Pica,
\textit{Topological susceptibility in SU(3) gauge theory},
Phys. Rev. Lett. 94 (2005) 032003

\bibitem{ChiInd}
M. L\"uscher,
\textit{Topological effects in QCD and the
problem of short distance singularities},
Phys. Lett. B593 (2004) 296

\bibitem{ChiProj}
L. Giusti, M. L\"uscher,
\textit{Chiral symmetry breaking and the Banks--Casher relation
in lattice QCD with Wilson quarks},
JHEP 0903 (2009) 013

\bibitem{ChiWI}
L. Giusti, G. C. Rossi, M. Testa, G. Veneziano,
\textit{The $U_A(1)$ problem on the lattice with Ginsparg--Wilson fermions},
Nucl. Phys. B628 (2002) 234

\bibitem{ChiWII}
L. Giusti, G. C. Rossi, M. Testa,
\textit{Topological susceptibility in full QCD with Ginsparg--Wilson fermions},
Phys. Lett. B587 (2004) 157

\end{thebibliography}



